\documentclass[carta,final]{lafrox}
\usepackage{simulacion}
\input{datos-proponente}
\usepackage{needspace}% keep headings with their content
\let\SIMoldsection\section
\renewcommand{\section}{\needspace{18\baselineskip}\SIMoldsection}
\begin{document}
\SIMtitulo{Estados financieros auditados}{LafroX SpA --- ejercicios 2023, 2024 y 2025}

\noindent\small\textit{Estados financieros simulados con fines académicos (Art.~34 y Art.~39°, numeral~2, de las Bases Administrativas). Cifras y auditor ficticios; sin valor contable ni legal. Cifras en millones de pesos chilenos (CLP MM), al 31 de diciembre de cada año.}\normalsize

\section*{Informe del auditor independiente}
\noindent A los señores accionistas de \LFXempresa:

\noindent \textbf{Opinión.} Hemos auditado los estados financieros de \LFXempresa, que comprenden los estados de situación financiera al 31 de diciembre de 2025, 2024 y 2023 y los correspondientes estados de resultados integrales y de flujos de efectivo por los años terminados en esas fechas, y las notas a los estados financieros. En nuestra opinión, los estados financieros presentan razonablemente, en todos sus aspectos significativos, la situación financiera de la Sociedad y los resultados de sus operaciones y flujos de efectivo, de acuerdo con las Normas Internacionales de Información Financiera. \textbf{Nuestra opinión es sin salvedades.}

\noindent \textbf{Fundamento de la opinión.} Efectuamos nuestras auditorías de acuerdo con las normas de auditoría generalmente aceptadas en Chile. Somos independientes de la Sociedad y hemos cumplido las demás responsabilidades éticas pertinentes. Consideramos que la evidencia de auditoría obtenida es suficiente y apropiada para fundamentar nuestra opinión.

\noindent \textbf{Responsabilidad de la administración.} La administración es responsable de la preparación y presentación razonable de los estados financieros, del control interno necesario para ello y de evaluar la capacidad de la Sociedad para continuar como empresa en marcha.

\vspace{6pt}
\noindent Santiago, 27 de febrero de 2026

\vspace{24pt}
\begin{center}
  \rule{6.4cm}{0.5pt}\\[2pt]
  Auditora Ficticia María Muestra Ejemplo, Socia\\
  \SIMauditora\\
  {\small\textit{Firma, timbre y registro de auditor simulados}}
\end{center}

\clearpage
\section*{Estado de situación financiera}
\begin{tablalafrox}{Estado de situación financiera (CLP MM).}{tab:ef1}{Y R{2.0cm} R{2.0cm} R{2.0cm}}
  {\cab{Concepto} & \cab{2025} & \cab{2024} & \cab{2023}}
  \textbf{Activos corrientes} & & & \\
  Efectivo y equivalentes al efectivo & 820 & 640 & 520 \\
  Deudores comerciales y otras cuentas por cobrar & 1.900 & 1.640 & 1.400 \\
  Otros activos corrientes & 430 & 360 & 330 \\
  \textbf{Total activos corrientes} & \textbf{3.150} & \textbf{2.640} & \textbf{2.250} \\
  \midrule
  \textbf{Activos no corrientes} & & & \\
  Propiedades, planta y equipos & 1.250 & 1.150 & 1.060 \\
  Activos intangibles (desarrollo de software) & 350 & 330 & 290 \\
  Otros activos no corrientes & 200 & 170 & 150 \\
  \textbf{Total activos no corrientes} & \textbf{1.800} & \textbf{1.650} & \textbf{1.500} \\
  \midrule
  \textbf{Total activos} & \textbf{4.950} & \textbf{4.290} & \textbf{3.750} \\
\end{tablalafrox}

\needspace{26\baselineskip}
\begin{tablalafrox}{Pasivos y patrimonio (CLP MM).}{tab:ef2}{Y R{2.0cm} R{2.0cm} R{2.0cm}}
  {\cab{Concepto} & \cab{2025} & \cab{2024} & \cab{2023}}
  \textbf{Pasivos corrientes} & & & \\
  Proveedores y otras cuentas por pagar & 780 & 690 & 640 \\
  Obligaciones laborales y previsionales & 640 & 560 & 500 \\
  Obligaciones financieras corrientes & 330 & 300 & 260 \\
  \textbf{Total pasivos corrientes} & \textbf{1.750} & \textbf{1.550} & \textbf{1.400} \\
  \midrule
  \textbf{Pasivos no corrientes} & & & \\
  Obligaciones financieras no corrientes & 900 & 750 & 680 \\
  Otros pasivos no corrientes & 100 & 90 & 70 \\
  \textbf{Total pasivos no corrientes} & \textbf{1.000} & \textbf{840} & \textbf{750} \\
  \midrule
  \textbf{Patrimonio} & & & \\
  Capital emitido & 1.000 & 1.000 & 1.000 \\
  Resultados acumulados & 755 & 480 & 242 \\
  Resultado del ejercicio & 445 & 420 & 358 \\
  \textbf{Total patrimonio} & \textbf{2.200} & \textbf{1.900} & \textbf{1.600} \\
  \midrule
  \textbf{Total pasivos y patrimonio} & \textbf{4.950} & \textbf{4.290} & \textbf{3.750} \\
\end{tablalafrox}

\section*{Estado de resultados integrales}
\begin{tablalafrox}{Estado de resultados por función (CLP MM).}{tab:ef3}{Y R{2.0cm} R{2.0cm} R{2.0cm}}
  {\cab{Concepto} & \cab{2025} & \cab{2024} & \cab{2023}}
  Ingresos de actividades ordinarias & 6.800 & 5.900 & 5.200 \\
  Costo de ventas & (4.700) & (4.130) & (3.640) \\
  \textbf{Margen bruto} & \textbf{2.100} & \textbf{1.770} & \textbf{1.560} \\
  Gastos de administración y ventas & (1.400) & (1.115) & (1.000) \\
  \textbf{Resultado operacional} & \textbf{700} & \textbf{655} & \textbf{560} \\
  Resultado financiero neto & (90) & (80) & (70) \\
  Resultado antes de impuestos & 610 & 575 & 490 \\
  Impuesto a las ganancias (27 \%) & (165) & (155) & (132) \\
  \textbf{Utilidad del ejercicio} & \textbf{445} & \textbf{420} & \textbf{358} \\
\end{tablalafrox}

\section*{Estado de flujos de efectivo (resumen)}
\begin{tablalafrox}{Flujos de efectivo por actividad (CLP MM).}{tab:ef4}{Y R{2.0cm} R{2.0cm} R{2.0cm}}
  {\cab{Concepto} & \cab{2025} & \cab{2024} & \cab{2023}}
  Flujo de actividades de la operación & 720 & 610 & 540 \\
  Flujo de actividades de inversión & (380) & (340) & (310) \\
  Flujo de actividades de financiación & (160) & (150) & (140) \\
  \midrule
  \textbf{Incremento neto de efectivo} & \textbf{180} & \textbf{120} & \textbf{90} \\
  Efectivo al inicio del año & 640 & 520 & 430 \\
  \textbf{Efectivo al final del año} & \textbf{820} & \textbf{640} & \textbf{520} \\
\end{tablalafrox}

\section*{Notas a los estados financieros (resumen)}
\begin{description}
  \item[Nota 1. Entidad.] \LFXempresa{} es una sociedad por acciones constituida en 2012, con casa matriz en Valparaíso y sede central en Santiago. Su actividad es el desarrollo, integración y operación de sistemas de software y servicios de tecnologías de la información.
  \item[Nota 2. Bases de preparación.] Los estados financieros se preparan de acuerdo con las NIIF, bajo el criterio de devengado y en pesos chilenos, redondeados a millones.
  \item[Nota 3. Reconocimiento de ingresos.] Los ingresos por desarrollo se reconocen según el avance de los hitos contractuales; los de operación y soporte, mensualmente conforme al servicio prestado.
  \item[Nota 4. Propiedades, planta y equipos e intangibles.] Se valoran al costo menos depreciación y amortización acumuladas. El desarrollo de software capitalizable se amortiza en cuatro años.
  \item[Nota 5. Obligaciones financieras.] Corresponden a créditos bancarios en pesos y UF, sin cláusulas de aceleración incumplidas al cierre.
  \item[Nota 6. Dividendos.] En 2024 y 2025 se pagaron dividendos por 120 y 145 millones, respectivamente, con cargo a las utilidades de los ejercicios anteriores.
  \item[Nota 7. Hechos posteriores.] No existen hechos posteriores al cierre que afecten significativamente los estados financieros.
\end{description}

\section*{Indicadores financieros}
\begin{tablalafrox}{Indicadores derivados de los estados financieros.}{tab:ef5}{Y R{2.0cm} R{2.0cm} R{2.0cm}}
  {\cab{Indicador} & \cab{2025} & \cab{2024} & \cab{2023}}
  Liquidez corriente (activo corriente / pasivo corriente) & 1,80 & 1,70 & 1,61 \\
  Endeudamiento (pasivo total / patrimonio) & 1,25 & 1,26 & 1,34 \\
  Patrimonio (CLP MM) & 2.200 & 1.900 & 1.600 \\
  Ventas anuales (CLP MM) & 6.800 & 5.900 & 5.200 \\
  Margen operacional (\%) & 10,3 & 11,1 & 10,8 \\
  Rentabilidad sobre patrimonio (\%) & 20,2 & 22,1 & 22,4 \\
\end{tablalafrox}

\noindent El patrimonio es positivo en los tres ejercicios, conforme al requisito de solvencia financiera del Art.~34.1 de las Bases Administrativas.

\SIMfirma{\SIMfecha}
\end{document}
